\documentclass[a4paper,twoside]{article}

\usepackage{morewrites}

\usepackage[svgnames,dvipsnames,x11names]{xcolor}
\usepackage{graphicx}
\usepackage{texsmith-lists}
\usepackage[english]{babel}
\usepackage[colorlinks=true, linkcolor=NavyBlue, urlcolor=NavyBlue, citecolor=NavyBlue]{hyperref}
\usepackage[a4paper]{geometry}
\usepackage{ts-fonts}
\usepackage{amsmath}
\usepackage{float}
\usepackage{seqsplit}
\usepackage{ts-code}

\usepackage{lastpage}
\usepackage{refcount}
\makeatletter
\def\ps@plain{
\let\@oddhead\@empty
\let\@evenhead\@empty
\def\@oddfoot{
\hfil
\ifnum\getpagerefnumber{LastPage}>1\relax
\thepage
\fi
\hfil}
\let\@evenfoot\@oddfoot
}
\makeatother

\title{Box Drawing\\[0.5em]\large Tree diagrams inside a listing}
\author{TeXSmith}\date{}
\begin{document}
\maketitle
\section{Diagrams drawn in a listing}
A listing that draws a tree uses the Unicode \emph{Box Drawing} block, and a
progress bar the \emph{Block Elements} one. No text font carries either, so the
document takes them from a monospace fallback --- at the advance width of the
column they sit in, which is what keeps a tree's stems under one another.
\begin{tscode}[lang=text, stretch=0.5, engine=pygments]
\begin{Verbatim}[commandchars=\\\{\},breaklines, breakanywhere, commandchars=\\\{\}]
project
├── src
│   ├── main.c
│   └── util.c
└── tests
    └── main\PYZus{}test.c

┌───┬───┐
│ a │ b │
├───┼───┤
│ c │ d │
└───┴───┘

[████████░░░░] 66\PYZpc{}
\end{Verbatim}
\end{tscode}
A Greek letter in a listing comes from the same machinery: \tscodeinline{eps = ε} is the
epsilon of \tscodeinline{if (fabs(x) < ε)}, and it is a fallback away from the monospace
face too.
\begin{tscode}[lang=c, engine=pygments]
\begin{Verbatim}[commandchars=\\\{\},breaklines, breakanywhere, commandchars=\\\{\}]
\PY{k+kt}{double}\PY{+w}{ }\PY{n}{eps}\PY{+w}{ }\PY{o}{=}\PY{+w}{ }\PY{l+m+mf}{1e\PYZhy{}9}\PY{p}{;}\PY{+w}{ }\PY{c+cm}{/* ε */}
\end{Verbatim}
\end{tscode}
\end{document}
